\begin{abstract}
Mathematical reasoning remains a formidable hurdle for language models due to its intricate and highly structured demands. This work presents DeepSeekMath~7B, which extends the pre-training of DeepSeek-Coder-Base-v1.5 7B using a dataset of 120B math-focused tokens derived from Common Crawl, in combination with data spanning natural language and code. DeepSeekMath~7B attains a notable 51.7\% on the competition-grade MATH benchmark without the use of auxiliary toolkits or voting strategies, nearing the benchmark achievements of Gemini-Ultra and GPT-4. Leveraging self-consistency with 64 samples, DeepSeekMath~7B yields a result of 60.9\% on MATH. The enhancement in DeepSeekMath~’s mathematical reasoning is ascribed to two primary innovations: firstly, we capitalize on the rich resources present in open web data by implementing a carefully curated data filtering pipeline; secondly, we propose Group Relative Policy Optimization (GRPO)—a modification of Proximal Policy Optimization (PPO)—which strengthens mathematical reasoning skills and simultaneously improves PPO’s memory efficiency.
\end{abstract}

\section{Introduction}

The advent of large language models (LLM) has transformed the landscape of mathematical reasoning within artificial intelligence, driving progress on benchmarks for quantitative \citep{MATH} as well as geometry-based reasoning \citep{trinh2024solving}. In addition, these models have become valuable tools in supporting humans to tackle advanced mathematical tasks \citep{tao}. Nonetheless, the highest-performing systems, including GPT-4 \citep{gpt4} and Gemini-Ultra \citep{gemini}, are not made available to the public, and the suite of open-source models currently falls well short of matching their performance.

This paper presents DeepSeekMath, a specialized language model focused on mathematics that demonstrates superior performance compared to existing open-source models and closely matches GPT-4 on academic evaluation metrics. Central to this achievement is the development of the DeepSeekMath Corpus—a comprehensive, high-caliber pre-training dataset consisting of 120B math-specific tokens. To assemble this corpus, we utilized the Common Crawl (CC) dataset, implementing a fastText-based classifier \citep{joulin2016fasttext} for selection. The classifier was initially trained with positive samples sourced from OpenWebMath \citep{openwebmath} and a variety of unrelated web pages as negative samples. Once trained, the classifier was deployed to identify additional relevant data from CC, and these selections were further vetted and annotated by human reviewers. The classifier underwent retraining with this enriched dataset, further enhancing its precision. Evaluations reveal the effectiveness of our approach, with DeepSeekMath-Base 7B recording 64.2\% accuracy on GSM8K \citep{gsm8k} and 36.2\% on the high-difficulty MATH dataset \citep{MATH}, surpassing Minerva 540B \citep{minerva}. Notably, the multilingual aspect of the DeepSeekMath Corpus contributes to improved outcomes on Chinese mathematical assessment sets \citep{wei2023cmath,agieval}. We see our methodology for constructing mathematical datasets as a valuable foundation for the broader research community, anticipating substantial future progress in this domain.

DeepSeekMath-Base is built upon DeepSeek-Coder-Base-v1.5 7B \citep{deepseek-coder}, motivated by our observation that initializing from a code-centric model yields superior outcomes relative to beginning with a general large language model. Additionally, our experiments reveal that mathematical pre-training not only augments the model's proficiency on mathematical tasks, but also enhances its performance on broader benchmarks such as MMLU \citep{mmlu} and BBH \citep{bbh}, demonstrating gains in general reasoning abilities alongside mathematical competence.

Subsequent to pre-training, we perform mathematical instruction tuning on DeepSeekMath-Base by incorporating datasets featuring chain-of-thought reasoning \citep{cot}, program-of-thought reasoning \citep{pot,pal}, and tool-augmented problem solving \citep{tora}. The resulting model, DeepSeekMath-Instruct 7B, outperforms all other models in the 7B parameter class and approaches the performance levels of open-source instruction-tuned models with 70B parameters.

Moreover, we propose Group Relative Policy Optimization (GRPO), an RL algorithm derived from Proximal Policy Optimization (PPO) \citep{schulman2017proximal}. Distinct from PPO, GRPO eliminates the critic model and instead derives the baseline using aggregated group scores, thereby substantially decreasing computational costs. Leveraging a subset of English instruction tuning data, GRPO delivers marked improvements over the already strong DeepSeekMath-Instruct, achieving notable performance gains across both in-domain (GSM8K: 82.9\% $\rightarrow$ 88.2\%, MATH: 46.8\% $\rightarrow$ 51.7\%) and out-of-domain mathematical evaluations (e.g., CMATH: 84.6\% $\rightarrow$ 88.8\%) during RL-based fine-tuning. We also articulate a unified framework that encompasses several existing approaches, such as Rejection Sampling Fine-Tuning (RFT) \citep{yuan2023scaling}, Direct Preference Optimization (DPO) \citep{dpo}, PPO, and GRPO. Within this unifying perspective, we recognize that these methodologies can be interpreted as direct or streamlined forms of reinforcement learning. To better understand the critical aspects of this framework, we perform extensive analyses, examining dimensions such as online versus offline training, supervision on outcomes versus process, and single-step versus iterative RL, among others. Finally, we elucidate the mechanisms by which RL contributes to the enhancement of instruction-tuned models, and we outline possible future directions for developing even more efficient RL strategies based on our unified conceptual paradigm.

\subsection{Contributions}
This work introduces advances in scalable mathematical pre-training and delves into the investigation of reinforcement learning approaches.

\textbf{Math Pre-Training at Scale}

\begin{itemize}[topsep=0pt]
    \item Our study demonstrates that Common Crawl, an openly available data resource, holds substantial potential for mathematical applications. Leveraging a carefully engineered pipeline for data selection, we assemble the DeepSeekMath Corpus—an extensive, high-quality dataset encompassing 120B tokens from web sources explicitly filtered for mathematical relevance. This corpus notably exceeds the scale of the math datasets employed by Minerva \citep{minerva} by a factor of nearly 7, and surpasses OpenWebMath \citep{openwebmath} by a factor of 9.

    \item The DeepSeekMath-Base 7B model, after pre-training, attains results on par with those of Minerva 540B \citep{minerva}, suggesting that parameter count alone is not the principal determinant of mathematical reasoning proficiency. Our findings highlight that smaller models, when trained on sufficiently curated datasets, can deliver robust performance.

    \item We present insights gleaned from math pre-training trials: engaging in code pre-training before mathematical tasks boosts the model's ability to address mathematical problems, regardless of whether external tools are used. This observation partially addresses the longstanding query of whether pre-training on code bolsters reasoning; our results support the view that it does—at least within the mathematical reasoning domain.

    \item While the practice of utilizing arXiv papers for training is prevalent within mathematical research, our experiments indicate that such data do not yield significant gains on the range of mathematical benchmarks evaluated in this study.
\end{itemize}

\textbf{Exploration and Analysis of Reinforcement Learning}

\begin{itemize}[topsep=0pt]
    \item We present Group Relative Policy Optimization (GRPO), a novel reinforcement learning approach characterized by its efficiency and efficacy. In contrast to Proximal Policy Optimization (PPO), GRPO dispenses with the use of a critic model and instead relies on group-based score estimation to determine the baseline. This design choice substantially lowers the computational burden during training.
    \item Our empirical analysis reveals that GRPO yields substantial improvements in the instruction-tuned DeepSeekMath-Instruct model, utilizing only instruction-tuning datasets. Additionally, we observe notable gains in generalization to out-of-domain tasks during reinforcement learning.
    \item We offer a unified conceptual framework for situating various techniques—including RFT, DPO, PPO, and GRPO—within a single paradigm. Comprehensive experimental studies are performed, comparing aspects such as online versus offline training modalities, outcome-based versus process-based supervision, and single-turn versus iterative reinforcement learning, to rigorously analyze the core constituents of this framework.
    \item Leveraging this unified paradigm, we investigate the mechanisms that underpin the efficacy of reinforcement learning and delineate several promising avenues for advancing large language model (LLM) reinforcement learning techniques.
\end{itemize}

\subsection{Summary of Evaluations and Metrics}

\begin{itemize}[topsep=0pt]
    \item \textbf{English and Chinese Mathematical Reasoning}: We carry out in-depth evaluations of our models using a broad array of English and Chinese mathematical benchmarks spanning grade-school through collegiate levels. The English assessment set comprises GSM8K \citep{gsm8k}, MATH \citep{MATH}, SAT \citep{llemma}, OCW Courses \citep{minerva}, and MMLU-STEM \citep{mmlu}. For Chinese, we utilize MGSM-zh \citep{mgsm}, CMATH \citep{wei2023cmath}, Gaokao-MathCloze \citep{agieval}, and Gaokao-MathQA \citep{agieval}. Evaluations target both the generation of independent, fully-detailed textual solutions (without recourse to external tools) and the capability to employ Python for problem solving.

    On the English benchmarks, DeepSeekMath-Base exhibits competitive performance vis-à-vis the closed-source Minerva 540B \citep{minerva}, and notably outperforms all open-source base models (e.g., Mistral 7B \citep{mistral}, Llemma-34B \citep{llemma}), regardless of prior math pre-training, often by considerable margins. Importantly, on Chinese benchmarks, DeepSeekMath-Base achieves superior results, which may be attributed to our inclusion of high-quality multilingual (not English-exclusive) pre-training data, diverging from earlier efforts \citep{minerva,llemma}. Mathematical instruction tuning combined with reinforcement learning leads to DeepSeekMath-Instruct and DeepSeekMath-RL, which attain robust outcomes, notably surpassing the 50\% accuracy threshold on the challenging competition-grade MATH dataset—a milestone not previously achieved by any open-source model.
\end{itemize}

\item \textbf{Formal Mathematics}: We assess the capabilities of DeepSeekMath-Base by applying it to the informal-to-formal theorem proving task introduced in \citep{dsp_proof}, utilizing the miniF2F benchmark \citep{minif2f} with Isabelle \citep{isabelle} as the designated proof assistant. The results reveal that DeepSeekMath-Base achieves strong performance in few-shot autoformalization settings.

\item \textbf{Natural Language Understanding, Reasoning, and Code}: To holistically examine the model's proficiency in general understanding, reasoning, and programming, we test DeepSeekMath-Base on the Massive Multitask Language Understanding (MMLU) benchmark \citep{mmlu}, which features 57 diverse multiple-choice tasks; BIG-Bench Hard (BBH) \citep{bbh}, composed of 23 multi-step reasoning challenges; as well as HumanEval \citep{codex} and MBPP \citep{mbpp}, two established benchmarks for evaluating code-related models. Math-specific pre-training enhances both comprehension and logical reasoning capabilities.

\section{Math Pre-Training}

\subsection{Data Collection and Decontamination} Here, we describe the methodology for assembling the DeepSeekMath Corpus utilizing Common Crawl data. As illustrated in Figure \ref{fig:data_collect}, an iterative approach is deployed, systematically collecting a vast mathematical dataset from Common Crawl, beginning with a seed dataset (for instance, a small, curated set of mathematical texts). Notably, this procedure generalizes to other fields such as programming.

To initiate the pipeline, OpenWebMath \citep{openwebmath}—a repository of meticulously selected mathematical web documents—serves as the foundational seed corpus. With this seed, we train a fastText model \citep{joulin2016fasttext} aimed at retrieving additional mathematical web content akin to OpenWebMath. Concretely, 500,000 samples from the seed are used as positive examples and 500,000 randomly chosen Common Crawl web pages as negative counterparts. The open-source implementation\footnote{\url{https://fasttext.cc}} is employed for training, with vector dimension set at 256, a learning rate of 0.1, word n-grams of up to length 3, a minimum word occurrence threshold of 3, and training proceeding for 3 epochs. To reduce dataset size, the original Common Crawl collection undergoes URL-based deduplication and near-duplicate filtering, yielding 40 billion HTML pages. Subsequently, mathematical pages are extracted from this deduplicated set with the aid of the fastText classifier. The relevance scores assigned by the model are then used to prioritize and retain only the highest-scoring mathematical pages, effectively filtering out lower quality material. The retained corpus size is empirically validated by conducting pre-training runs with datasets composed of the top 40B, 80B, 120B, and 160B tokens, and for the initial iteration, we select the highest-scoring 40B tokens.

Following the initial round of data gathering, a significant number of mathematical web pages remain uncollected. This is primarily attributed to the limited diversity within the set of positive examples used to train the fastText model. To address this, we augment the seed dataset by sourcing additional mathematical websites, thereby enhancing the capacity of the fastText classifier. Concretely, we partition the entire Common Crawl dataset into mutually exclusive domains, where each domain consists of web pages with a shared base URL. For every domain, we compute the proportion of pages that were obtained during the initial pass. If a domain has over 10\% of its pages collected, it is flagged as likely math-related (for instance, \url{mathoverflow.net}). 

Next, we manually inspect and annotate URLs that are linked to mathematical content within these domains (e.g., \url{mathoverflow.net/questions}). Any previously uncollected web pages associated with these annotated URLs are subsequently incorporated into our seed corpus. This strategy permits us to increase the number of positive examples, which, in turn, yields a more robust fastText model capable of retrieving a greater volume of mathematical data in the following round. After completing four collection cycles, we acquire 35.5M mathematical web pages, comprising 120B tokens. By the fourth round, analysis reveals that approximately 98\% of the total pages were already captured in the preceding round, prompting us to halt further collection.

To prevent data contamination in evaluation benchmarks, we adopt the method of \cite{deepseek-coder} and exclude any web pages containing questions or answers sourced from English mathematical benchmarks (such as GSM8K \citep{gsm8k} and MATH \citep{MATH}) as well as from Chinese benchmarks (including CMATH \citep{wei2023cmath} and AGIEval \citep{agieval}). The removal protocol specifies that any text segment containing an exact 10-gram overlap with a substring from any benchmark is discarded from the mathematical training corpus. For benchmark strings shorter than 10 grams but with a minimum length of 3 grams, we filter via precise matching to identify and remove any contaminated content.

\subsection{Validating the Quality of the DeepSeekMath~Corpus}

We conduct pre-training experiments to assess the performance of the DeepSeekMath~Corpus relative to the latest released mathematics-oriented corpora:

\begin{itemize}[topsep=0pt]
    \item \textbf{MathPile} \citep{mathpile}: a composite corpus spanning multiple sources (8.9B tokens), drawn from textbooks, Wikipedia, ProofWiki, CommonCrawl, StackExchange, and arXiv, with the bulk of the data (over 85\%) derived from arXiv;
    \item \textbf{OpenWebMath} \citep{openwebmath}: a mathematical data subset of CommonCrawl, comprising 13.6B tokens;
    \item \textbf{Proof-Pile-2} \citep{llemma}: a math-specific corpus that incorporates OpenWebMath, AlgebraicStack (which contains 10.3B tokens of mathematical code), and arXiv manuscripts (28.0B tokens). For Proof-Pile-2 evaluations, we follow the arXiv:Web:Code sampling ratio of 2:4:1 as described in \cite{llemma}.
\end{itemize}

\subsubsection{Training Setting}
\label{sec:quality-policy}
We conduct specialized mathematical training on a general-purpose language model comprising 1.3B parameters, employing the same architecture as DeepSeek LLMs \citep{deepseek-llm}—hereafter referred to as DeepSeek-LLM 1.3B. Each mathematical dataset is used to independently train a model for a total of 150B tokens. The entirety of our experimental procedures leverages the resource-efficient and lightweight HAI-LLM \citep{haillm} training infrastructure. In line with the established DeepSeek LLMs methodology, we utilize the AdamW optimizer \citep{adamW} with $\beta_1=0.9$, $\beta_2=0.95$, and $\mathrm{weight\_decay}=0.1$. A multi-stage learning rate strategy is adopted: the rate increases to its maximum after 2,000 warmup steps, then decreases to 31.6\% of this maximum by the time 80\% of total training steps have elapsed, and further declines to 10.0\% of the peak by 90\% progress. The upper bound for the learning rate is set at 5.3e-4. Training proceeds with a batch size accommodating 4 million tokens and a sequence context window of 4K tokens.

\subsubsection{Evaluation Results}

\textbf{The DeepSeekMath~Corpus is distinguished by its high quality, broad multilingual scope, and unparalleled scale.}

\begin{itemize}[topsep=0pt]
    \item \textbf{High-quality}:
    We measure downstream performance on eight mathematics-focused benchmarks utilizing few-shot chain-of-thought prompting \cite{cot}. The results in Table \ref{tab:corpora_comparison} indicate that models trained on the DeepSeekMath~Corpus distinctly outperform their counterparts. Furthermore, Figure \ref{fig:corpora_comparisons} illustrates that after processing 50B tokens—the equivalent of one complete pass over Proof-Pile-2—the DeepSeekMath Corpus-trained model exhibits superior capabilities compared to Proof-Pile-2, suggesting a higher average data quality for DeepSeekMath Corpus.
    \item \textbf{Multilingual}:
    The DeepSeekMath~Corpus incorporates mathematical data across several languages, with English and Chinese being the most prevalent. Evidence in Table \ref{tab:corpora_comparison} demonstrates that training with DeepSeekMath~Corpus markedly enhances mathematical reasoning in both English and Chinese. By comparison, other mathematical datasets, which tend to be primarily in English, yield limited or even negative impact on Chinese mathematical tasks.
    \item \textbf{Large-scale}:
    In terms of volume, the DeepSeekMath~Corpus substantially surpasses other mathematical corpora. Figure \ref{fig:corpora_comparisons} reveals that DeepSeek-LLM 1.3B, when trained on this corpus, achieves more pronounced and sustained learning gains. Meanwhile, existing baseline corpora, being considerably smaller, are subjected to repeated epochs during training, with model performance stagnating after an initial rise.
\end{itemize}

\subsection{Training and Evaluating DeepSeekMath-Base 7B}

This section details DeepSeekMath-Base 7B, a foundational model designed to demonstrate advanced reasoning, with a particular focus on mathematical domains. The initial weights are derived from DeepSeek-Coder-Base-v1.5 7B \citep{deepseek-coder}, followed by pretraining over 500B tokens. The data composition consists of 56\% DeepSeekMath Corpus, 4\% AlgebraicStack, 10\% arXiv articles, 20\% Github code, and the final 10\% from Common Crawl, incorporating both English and Chinese natural language samples. The training regimen aligns primarily with Section \ref{sec:quality-policy}, with modifications setting the peak learning rate at 4.2e-4 and a batch size of 10M tokens.

A thorough evaluation of DeepSeekMath-Base 7B's mathematical proficiency is performed, examining its capability to independently construct mathematical solutions, utilize external tools for problem-solving, and engage in formal theorem proving. Additionally, the model's broader abilities—including comprehension of natural language, logical reasoning, and programming performance—are systematically reported.

\paragraph{Mathematical Problem Solving with Step-by-Step Reasoning}

The mathematical problem-solving effectiveness of DeepSeekMath-Base is assessed using few-shot chain-of-thought prompting \citep{cot} across eight benchmarks in both English and Chinese. These tasks span quantitative reasoning challenges (such as GSM8K \citep{gsm8k}, MATH \citep{MATH}, and CMATH \citep{wei2023cmath}) and multiple-choice formats (like MMLU-STEM \citep{mmlu} and Gaokao-MathQA \citep{agieval}), thereby covering a broad spectrum of mathematics topics, ranging from elementary mathematics to college-level problems.

As detailed in Table \ref{tab:base_math_cot}, DeepSeekMath-Base 7B achieves the highest scores across all eight evaluated benchmarks when compared with open-source base models, including both the widely adopted Mistral 7B \citep{mistral} and the more recent Llemma 34B \citep{llemma}, which was specially trained on Proof-Pile-2 \citep{llemma} for mathematics. In particular, DeepSeekMath-Base sets itself apart on the challenging MATH competition dataset, exceeding the best-performing open-source alternatives by a margin of over 10\% in absolute terms, and even surpassing Minerva 540B \citep{minerva}—a proprietary model based on PaLM \citep{palm} and further refined on mathematical corpora—with only a fraction of the parameter count.

\paragraph{Mathematical Problem Solving with Tool Use} For the assessment of program-based mathematical reasoning, we apply few-shot program-of-thought prompting \citep{pot,pal} on GSM8K and MATH. Here, models are guided to respond by constructing Python code, with the use of libraries like \textit{math} and \textit{sympy} permitted for complex calculations. The correctness of each solution is determined by the output of the corresponding program. According to Table \ref{tab:base_math_pot_proof}, DeepSeekMath-Base 7B surpasses Llemma 34B, which previously held the best results among open-source base models.

\paragraph{Formal Mathematics}

The automation of formal proof generation provides important benefits for proof reliability, correctness, and productivity, and has attracted considerable interest recently. We evaluate DeepSeekMath-Base 7B on the informal-to-formal proof generation task outlined in \citep{dsp_proof}, where the objective is to produce a formal proof given an informal problem statement, its formal version, and an accompanying informal argument. For this, we make use of miniF2F \citep{minif2f}, a formal mathematics benchmark focused on Olympiad-level problems, generating proofs in Isabelle using few-shot examples. Consistent with the procedure in \cite{dsp_proof}, models first generate high-level proof sketches, which are subsequently completed using the automated prover Sledgehammer \citep{sledgehammer}. Table \ref{tab:base_math_pot_proof} highlights the strong capabilities of DeepSeekMath-Base 7B in automatic proof formalization.

\paragraph{Natural Language Understanding, Reasoning, and Code}

We test model abilities for language comprehension using MMLU \citep{mmlu}, reasoning using BBH \citep{bbh}, and programming skills with HumanEval \citep{codex} and MBPP \citep{mbpp}. The results in Table \ref{tab:understanding_reasoning_code} show that DeepSeekMath-Base 7B makes substantial advances on MMLU and BBH relative to its predecessor, DeepSeek-Coder-Base-v1.5 \citep{deepseek-coder}, demonstrating the significant gains that specialized mathematical training offers for both language understanding and logical reasoning. Furthermore, the retention of code token exposure throughout continued training enables DeepSeekMath-Base 7B to maintain the programming abilities of DeepSeek-Coder-Base-v1.5 on coding evaluations. Overall, DeepSeekMath-Base 7B provides a marked improvement over the Mistral 7B \citep{mistral} general model on all three evaluated reasoning and coding tasks.

\section{Supervised Fine-Tuning}

\subsection{SFT Data Curation}

We assemble a dataset for mathematical instruction tuning that incorporates problems in both English and Chinese, spanning a wide array of mathematical domains and complexity. Each problem is paired with corresponding solutions presented in chain-of-thought (CoT) \citep{cot}, program-of-thought (PoT) \citep{pot,pal}, or tool-integrated reasoning formats \citep{tora}. Our final training corpus consists of 776K instances.

\begin{itemize}[topsep=0pt]
    \item \textbf{English mathematical datasets}: Our dataset contains GSM8K and MATH problems annotated with tool-integrated solutions. We also include selected examples from MathInstruct \citep{MathInstruct}, as well as the training split from Lila-OOD \citep{lila}, both featuring solutions generated via CoT or PoT techniques. This English portion encompasses mathematical subfields including algebra, probability, number theory, calculus, and geometry.
    \item \textbf{Chinese mathematical datasets}: The Chinese collection comprises K-12 math problems across 76 topics such as linear equations, each annotated with both CoT-style and tool-integrated solution formats.
\end{itemize}

\subsection{Training and Evaluating DeepSeekMath-Instruct 7B}

Here, we detail the instruction tuning process applied to DeepSeekMath-Instruct 7B, using DeepSeekMath-Base as the foundation. Training examples are concatenated at random until the input sequence length reaches up to 4K tokens. The training regimen consists of 500 steps, utilizing a batch size of 256, and a fixed learning rate of 5e-5.

For model evaluation, we assess mathematical capabilities in scenarios with and without the assistance of external tools, using four quantitative reasoning benchmarks available in both English and Chinese. Comparative benchmarking includes top-performing contemporaneous models:

\begin{itemize}[topsep=0pt]
    \item \textbf{Closed-source models} considered: (1) various GPT models, with GPT-4 \citep{gpt4} and GPT-4 Code Interpreter~\footnote{\url{https://openai.com/blog/chatgpt-plugins\#\#code-interpreter}} noted as state-of-the-art, (2) Gemini Ultra and Pro \citep{gemini}, (3) Inflection-2 \citep{inflection-2}, (4) Grok-1~\footnote{\url{https://x.ai/model-card}}, and prominent Chinese models such as (5) Baichuan-3~\footnote{\url{https://www.baichuan-ai.com}}, (6) the newly released GLM-4~\footnote{\url{https://open.bigmodel.cn/dev/api\#glm-4}} from the GLM series \citep{glm}.
\end{itemize}

These models are designed for general applications and typically have been subjected to alignment processes.

\item \textbf{Open-source models} consist of both general-purpose and math-specialized systems: (1) DeepSeek-LLM-Chat 67B \citep{deepseek-llm}, (2) Qwen 72B \citep{qwen}, (3) SeaLLM-v2 7B \citep{seallm}, and (4) ChatGLM3 6B \citep{chatglm3} are general models. Several others focus on mathematical tasks, such as (5) InternLM2-Math 20B~\footnote{\url{https://github.com/InternLM/InternLM-Math}}, built upon InternLM2 and further trained on mathematical data before undergoing instruction tuning, (6) Math-Shepherd-Mistral 7B, which employs PPO training \citep{schulman2017proximal} with a process-supervised reward mechanism applied to Mistral 7B \citep{mistral}, (7) the WizardMath family \citep{wizardmath}, which augments mathematical reasoning in Mistral 7B and Llama-2 70B \citep{llama2} through evolve-instruct (AI-generated instructions for instruction tuning) and PPO, using mainly problems from GSM8K and MATH, (8) MetaMath 70B \citep{metamath}, which refines Llama-2 70B via fine-tuning on extended GSM8K and MATH datasets, (9) ToRA 34B \cite{tora}, a CodeLlama 34B variant adjusted for tool-assisted mathematical reasoning, and (10) MAmmoTH 70B \citep{MathInstruct}, an instruction-tuned version of Llama-2 70B on MathInstruct.

As presented in Table \ref{tab:sft_rl_math}, in settings where external tool usage is prohibited, DeepSeekMath-Instruct 7B displays exceptional step-wise reasoning skills. Specifically, for the challenging MATH dataset, this model exceeds all open-source peers and most proprietary alternatives (such as Inflection-2 and Gemini Pro) by a minimum of 9\% absolute margin. This advantage persists against considerably larger models (e.g., Qwen 72B) and those trained with intensive math-centered reinforcement learning techniques (e.g., WizardMath-v1.1 7B). Although DeepSeekMath-Instruct attains performance on par with leading Chinese proprietary systems like GLM-4 and Baichuan-3 for MATH, it does not surpass GPT-4 or Gemini Ultra.

When models are permitted to combine natural language reasoning with code-based tools, DeepSeekMath-Instruct 7B achieves nearly 60\% accuracy on MATH, outperforming all other open-source models. For other datasets, its results are comparable to those of DeepSeek-LLM-Chat 67B, which is tenfold larger and was previously state-of-the-art.

\section{Reinforcement Learning}

\subsection{Group Relative Policy Optimization} Reinforcement learning (RL) continues to demonstrate its efficacy for boosting LLMs’ mathematical reasoning following Supervised Fine-Tuning (SFT) \citep{wang2023math,wizardmath}. This section details our novel and efficient RL approach, Group Relative Policy Optimization (GRPO).

\subsubsection{From PPO to GRPO} Proximal Policy Optimization (PPO) \citep{schulman2017proximal} stands out as a prominent actor-critic RL method and is extensively adopted in the RL fine-tuning of LLMs \citep{ouyang2022training}. It aims to optimize model policies through maximization of the following surrogate loss:

\begin{equation}
\footnotesize
    \mathcal{J}_{PPO}(\theta) = \mathbb{E}{[q \sim P(Q), o \sim \pi_{\theta_{old}}(O|q)]} \frac{1}{|o|} \sum_{t=1}^{|o|} \min \left[ \frac{\pi_\theta(o_{t} | q, o_{<t})}{\pi_{\theta_{old}}(o_{t} | q, o_{<t})} A_{t}, \text{clip} \left( \frac{\pi_\theta(o_{t} | q, o_{<t})}{\pi_{\theta_{old}}(o_{t} | q, o_{<t})}, 1 - \varepsilon, 1 + \varepsilon \right)  A_{t} \right] ,
\end{equation}

Here, $\pi_{\theta}$ denotes the current policy, while $\pi_{\theta_{old}}$ refers to the previous version of the policy model. The variables $q$ and $o$ represent questions drawn from the dataset and their associated outputs, with the latter sampled according to the old policy $\pi_{\theta_{old}}$. The hyper-parameter $\varepsilon$ is employed as part of the clipping mechanism in PPO to enhance the stability of optimization. The advantage $A_t$ is estimated using Generalized Advantage Estimation (GAE) \citep{gae}, which leverages both the reward sequence $\{r_{\ge t}\}$ and a parameterized value function $V_{\psi}$. In PPO, this setup necessitates the concurrent training of a value function with the policy model. To reduce excessive optimization towards the reward model, a per-token Kullback-Leibler (KL) penalty is generally incorporated into the reward at each token step, utilizing a reference model as outlined by \citep{ouyang2022training}:

\begin{equation}
    r_{t} = r_\varphi(q, o_{\le t}) - \beta \log\frac{\pi_{\theta}(o_{t}|q, o_{<t})}{\pi_{ref}(o_{t}|q, o_{<t})},
\label{eq:PPO-reward}
\end{equation}

where $r_\varphi$ indicates the reward function, $\pi_{ref}$ is the reference policy—typically instantiated as the original SFT model—and $\beta$ controls the strength of the KL term.

Because the value function in PPO is typically architecturally similar and almost as large as the policy model, it significantly increases memory requirements and computational complexity. Moreover, in reinforcement learning, this value function serves as a baseline to help reduce the variance of the advantage calculation. In practice, especially within large language model settings, only the final token usually receives a reward from the reward model, thereby complicating efforts to train a value function that is precise across all time steps. To address this issue, we introduce Group Relative Policy Optimization (GRPO), as depicted in Figure \ref{fig:grpo}. GRPO dispenses with the need for a separately parameterized value function, unlike PPO. Instead, it uses the mean reward from a set of output samples—each produced in response to the same question—as a baseline for normalization. Specifically, for each question $q$, GRPO generates a set of outputs $\{o_1, o_2, \cdots, o_G\}$ sampled via $\pi_{\theta_{old}}$, and optimizes the policy by maximizing the following expected objective:

\begin{equation}
\footnotesize
\begin{split}
    \mathcal{J}_{GRPO}(\theta) &= \mathbb{E}{[q \sim P(Q), \{o_i\}_{i=1}^G \sim \pi_{\theta_{old}}(O|q)]}  \\
    & \frac{1}{G}\sum_{i=1}^G\frac{1}{|o_i|} \sum_{t=1}^{|o_i|} \left\{ \min \left[ \frac{\pi_\theta(o_{i,t} | q, o_{i,<t})}{\pi_{\theta_{old}}(o_{i,t} | q, o_{i,<t})} \hat{A}_{i,t}, \text{clip} \left( \frac{\pi_\theta(o_{i,t} | q, o_{i,<t})}{\pi_{\theta_{old}}(o_{i,t} | q, o_{i,<t})}, 1 - \varepsilon, 1 + \varepsilon \right)  \hat{A}_{i,t} \right] - \beta \mathbb{D}_{KL}\left[\pi_{\theta} || \pi_{ref}\right]\right\} ,
\end{split}
\label{eq:GRPO-obj}
\end{equation}

where $\varepsilon$ and $\beta$ act as hyper-parameters, and $\hat{A}_{i,t}$ denotes the group-relative advantage that is derived solely based on relative rewards among the sampled outputs for a given question—details of which will be provided in the subsequent sections. By using group-wise comparisons for advantage computation, GRPO aligns closely with the reward models’ intrinsic preference-based structure, since these models are often trained on data comprising pairs or sets of outputs in response to a shared prompt. It should also be highlighted that GRPO directly imposes the KL regularization between the updated and reference policies on the loss function, rather than through reward augmentation, which streamlines the estimation of $\hat

\begin{algorithm}[t]
  \small
  \caption{Iterative Group Relative Policy Optimization}
  \textbf{Input} initial policy model $\pi_{\theta_{\text{init}}}$; reward models $r_\varphi$; task prompts $\mathcal{D}$; 
  hyperparameters $\varepsilon$, $\beta$, $\mu$
  \begin{algorithmic}[1]
    \State Set policy model $\pi_\theta$ to $\pi_{\theta_{\text{init}}}$
    \For{iteration = 1, \dots, I}
      \State Assign reference model $\pi_{ref} \leftarrow \pi_{\theta}$
      \For{step = 1, \dots, M}
      \State Draw a minibatch $\mathcal{D}_b$ from $\mathcal{D}$
      \State Update prior policy model $\pi_{\theta_{old}} \leftarrow \pi_{\theta}$ 
      \State For each question $q \in \mathcal{D}_b$, generate $G$ candidate responses $\{o_i\}_{i=1}^G \sim \pi_{\theta_{old}} (\cdot \mid q)$
      \State Evaluate each candidate $o_i$ using $r_{\varphi}$ to obtain rewards $\{r_i\}_{i=1}^{G}$
      \State Estimate group-relative advantages $\hat{A}_{i,t}$ at each position $t$ for each $o_i$
      \For{GRPO iteration = 1, \dots, $\mu$}
        \State Update policy model $\pi_{\theta}$ by maximizing the GRPO loss in Equation \ref{eq:GC-GRPO}
      \EndFor
    \EndFor \State Continuously retrain $r_\varphi$ using a replay-based scheme.
    \EndFor 
  \end{algorithmic}
  \textbf{Output} $\pi_\theta$
  \label{alg:iter-grpo}
\end{algorithm}

\subsubsection{Outcome Supervision RL with GRPO} Given a question $q$, a set of outputs $\{o_1, o_2, \cdots, o_G\}$ are sampled from $\pi_{\theta_{old}}$. Each output is then assigned a reward using a reward model, yielding $\mathbf{r}=\{r_1, r_2, \cdots, r_G\}$. These rewards are normalized by subtracting their mean and dividing by their standard deviation within the group. With outcome supervision, the final normalized reward for output $o_i$ is used as the advantage for every token in $o_i$, namely $\hat{A}_{i, t} = \widetilde{r}_i = \frac{r_i- {\rm mean}(\mathbf{r})}{{\rm std}(\mathbf{r})}$. The policy is then optimized by maximizing the target specified in equation (\ref{eq:GRPO-obj}).

\subsubsection{Process Supervision RL with GRPO} Outcome supervision provides feedback solely at the end of the output, which may not provide sufficient guidance for challenging mathematical reasoning. Inspired by \cite{wang2023math}, we also implement process supervision, delivering a reward after each intermediate reasoning stage. More precisely, for a given question $q$ and a collection of generated outputs $\{o_1, o_2, \cdots, o_G\}$, a process reward model is employed to assign scores to each step in the solutions, resulting in reward sequences $\mathbf{R} = \{ \{r_1^{index(1)},\cdots,r_1^{index(K_1)}\}, \cdots,  \{r_G^{index(1)},\cdots,r_G^{index(K_G)}\} \}$, with $index(j)$ denoting the position of the final token in step $j$, and $K_i$ the number of steps in $o_i$. These scores are normalized by computing $\widetilde{r}_i^{index(j)} = \frac{r_i^{index(j)} - {\rm mean(\mathbf{R})}}{{\rm std(\mathbf{R})}}$. The advantage for each token $t$ in $o_i$ is then determined as the sum of all normalized step rewards at positions no earlier than $t$: $\hat{A}_{i, t} = \sum_{index(j) \ge t} \widetilde{r}_i^{index(j)}$. Policy optimization then proceeds by maximizing the objective defined by equation (\ref{eq:GRPO-obj}).

\subsubsection{Iterative RL with GRPO} During reinforcement learning, as training advances, the previous reward model might become inadequate for supervising the evolving policy model. To address this, we investigate iterative RL with GRPO. Algorithm \ref{alg:iter-grpo} illustrates this approach: in each iteration, new training datasets for the reward model are synthesized by sampling outputs from the policy model, and the reward model is further trained via a replay strategy that incorporates 10\% of prior samples. Subsequently, the updated policy model is set as the new reference, and further policy model training is conducted under supervision from the updated reward model.

\subsection{Training and Evaluating DeepSeekMath-RL}

Our reinforcement learning procedures are based on DeepSeekMath-Instruct 7B. The RL training data comprise chain-of-thought formatted items derived from GSM8K and MATH within the SFT dataset, amounting to approximately 144K questions. Other SFT questions are omitted, allowing us to specifically assess the effect of RL on evaluation benchmarks where no corresponding training examples are present during RL. The preparation of the reward model’s training set follows the methodology of \citep{wang2023math}. The initial reward model is trained using DeepSeekMath-Base 7B with a learning rate of 2e-5. For GRPO training, we use a policy model learning rate of 1e-6 and set the KL penalty coefficient to 0.04. For each item, $64$ responses are generated. The maximum output sequence length is 1024, and the batch size for training is likewise set to 1024. Each policy model receives only one update per exploration phase. Evaluations of DeepSeekMath-RL 7B are conducted with the same benchmarks as DeepSeekMath-Instruct 7B. For DeepSeekMath-RL 7B, GSM8K and MATH with chain-of-thought solutions are considered in-domain benchmarks, while all other evaluation benchmarks are classified as out-of-domain.

Table \ref{tab:sft_rl_math} presents comparative results for both open-source and proprietary models employing chain-of-thought and tool-assisted reasoning across English and Chinese evaluation datasets. Our analysis reveals: 1) DeepSeekMath-RL 7B achieves chain-of-thought accuracies of 88.2\% on GSM8K and 51.7\% on MATH. These scores not only exceed all open-source competitors within the 7B--70B parameter range, but also outperform most closed-source alternatives. 2) Importantly, DeepSeekMath-RL 7B's training regimen is limited to chain-of-thought-style instruction tuning data drawn solely from GSM8K and MATH, with DeepSeekMath-Instruct 7B serving as the initialization point. Despite its narrow training corpus, DeepSeekMath-RL 7B delivers superior results compared to DeepSeekMath-Instruct 7B on every metric, underscoring the efficacy of the applied reinforcement learning methods.

\section{Discussion}
This section details insights gleaned from our pre-training and reinforcement learning investigations.

\subsection{Lessons Learnt in Pre-Training}

We begin by recounting our observations from pre-training. Unless otherwise noted, all training configurations mirror those described in Section \ref{sec:quality-policy}. For clarity, in this context, the DeepSeekMath Corpus refers specifically to the 89B-token version curated in the data collection’s second round.

\subsubsection{Code Training Benefits Mathematical Reasoning}

There is a prevalent, though previously unproven, assumption that code-related training enhances reasoning abilities. Our findings contribute evidence supporting this notion in mathematics: incorporating code into training boosts the capacity for mathematical reasoning, independent of tool utilization.

To assess the impact of code training on mathematical reasoning skills, we performed comparative studies involving both two-stage and one-stage training protocols:

\textbf{Two-Stage Training}

\begin{itemize}[topsep=0pt]
    \item \textbf{Code Training for 400B Tokens $\rightarrow$ Math Training for 150B Tokens}: The DeepSeek-LLM 1.3B model is first trained on a dataset consisting of 400B code tokens, after which it undergoes further training on 150B math-specific tokens;
    \item \textbf{General Training for 400B Tokens $\rightarrow$ Math Training for 150B Tokens}: For comparison, we conduct a control experiment where, instead of code tokens, we utilize 400B general corpus tokens (sourced from DeepSeek-AI’s large-scale data) in the first phase to analyze whether code-centric pre-training offers distinct advantages over generic data for improving mathematical reasoning.
\end{itemize}

\textbf{One-Stage Training}

\begin{itemize}[topsep=0pt]
    \item \textbf{Math Training for 150B Tokens}: In this setup, DeepSeek-LLM 1.3B is exposed solely to 150B math tokens during training;
    \item \textbf{Training on a mixture of 400B Code Tokens and 150B Math Tokens}: Since sequential math training post code-pretraining adversely affects code task performance, we explore a joint-training scenario, interleaving 400B code tokens and 150B math tokens, to determine whether concurrent exposure supports both improved mathematical reasoning and retention of coding proficiency.
\end{itemize}

\paragraph{Results}
Tables \ref{tab:code-math} and \ref{tab:code-math-general-eval} report the outcomes on downstream evaluation tasks for the different training regimes.

Pre-training on code data is beneficial for program-based mathematical problem solving in both sequential (two-stage) and joint (one-stage) training paradigms. According to Table \ref{tab:code-math}, in the two-stage approach, code pre-training substantially raises performance on GSM8K and MATH tasks utilizing Python. Subsequent math-focused fine-tuning further amplifies this effect. Notably, with the one-stage joint approach, integrating code and math tokens reduces the risk of catastrophic forgetting, which otherwise undermines code ability in sequential setups, and maintains gains in both coding (see Table \ref{tab:code-math-general-eval}) and program-assisted mathematical reasoning (see Table \ref{tab:code-math}).

Code-centric training also enhances direct mathematical reasoning without code execution. In the two-stage process, the initial exposure to code tokens yields moderate gains, and it facilitates improved effectiveness of follow-up math training, producing the strongest results overall. Nevertheless, blending code and math tokens in a single-stage curriculum is less effective for math reasoning without code assistance. This may be attributed to the constrained capacity of DeepSeek-LLM 1.3B, which likely hampers its ability to concurrently learn and retain both coding and mathematical skills at this model scale.

\subsubsection{ArXiv Papers Appear Not to Enhance Mathematical Reasoning}

ArXiv papers are frequently incorporated as a substantial part of mathematical pre-training datasets \citep{minerva,gpt-f,llemma,mathpile}. Despite their prevalence, the specific influence of arXiv content on mathematical reasoning capabilities remains insufficiently studied. Surprisingly, our empirical results indicate that arXiv papers do not contribute positively to mathematical reasoning ability. We assess this by training models of varying scale, including DeepSeek-LLM 1.3B and DeepSeek-Coder-Base-v1.5 7B \citep{deepseek-coder}, on arXiv collections subjected to diverse preprocessing approaches:

\begin{itemize}[topsep=0pt]
    \item \textbf{MathPile} \citep{mathpile}: a corpus containing 8.9B tokens, processed through rigorous cleaning and filtering heuristics, with over 85\% of its content originating from scientific arXiv documents;
    \item \textbf{ArXiv-RedPajama} \citep{redpajama}: comprising the full set of arXiv LaTeX source files, stripped of preambles, comments, custom macros, and reference lists, summing to 28.0B tokens.
\end{itemize}

In our evaluation, DeepSeek-LLM 1.3B is trained on each arXiv dataset for 150B tokens, and DeepSeek-Coder-Base-v1.5 7B for 40B tokens. Contrary to common expectations, we observe that exclusive training on arXiv data fails to produce meaningful gains, and occasionally even leads to regression, in various mathematical evaluation benchmarks covering a spectrum of difficulties. The benchmarks examined include quantitative reasoning datasets such as GSM8K and MATH (refer to Table \ref{tab:arxiv-cot}), multiple-choice STEM evaluations like MMLU-STEM (Table \ref{tab:arxiv-cot}), and benchmarks for formal mathematics including miniF2F (Table \ref{tab:arxiv_atp}).

Nevertheless, these findings are not definitive and should be interpreted with caution. Specifically, we have not examined:

\begin{itemize}[topsep=0pt]
    \item The influence of arXiv-derived tokens on specialized mathematical tasks not assessed in our current scope, such as the informalization process, which involves translating formal mathematical statements or proofs into informal language;
    \item The effect that arXiv data might have when it is integrated with additional data sources;
    \item The potential advantages that might emerge from using arXiv papers with models at greater parameter scales.
\end{itemize}

Consequently, we highlight the need for additional, more comprehensive investigation on this topic in future work.

\subsection{Insights of Reinforcement Learning}
\subsubsection{Toward a Unified Training Framework} Here, we introduce a generalized framework for comparing a variety of training strategies—including SFT, RFT, DPO, PPO, and GRPO—and conduct experiments probing the components of this unified formulation. The parameter gradient $\theta$ for any of these training procedures can generally be represented as:

\begin{equation}
    \nabla_{\theta}\mathcal{J}_{\textcolor{red}{\mathcal{A}}}(\theta) = \mathbb{E}[\underbrace{(q,o) \sim \textcolor{red}{\mathcal{D}}}_{Data \ Source}]\left( \frac{1}{|o|} \sum_{t=1}^{|o|}  \underbrace{GC_{{\mathcal{A}}}(q, o, t, \textcolor{red}{\pi_{{rf}}})}_{Gradient \ Coefficient}  \nabla_{\theta}\log \pi_{\theta}(o_t | q, o_{<t})\right).
\label{eq:objective}
\end{equation}

The formulation above highlights three principal elements: 1) \textit{Data Source $\mathcal{D}$}, specifying the origin of the training samples; 2) \textit{Reward Function $\pi_{{rf}}$}, providing the feedback signal to guide learning; and 3) \textit{Algorithm $\mathcal{A}$}, which processes both training inputs and reward feedback into a gradient coefficient $GC$ that modulates the extent of learning—either through penalization or reinforcement. We organize several common approaches within this general framework:

\begin{itemize}
    \item  \textbf{Supervised Fine-tuning (SFT)}: This method involves additional training of the pretrained model using curated SFT datasets constructed by human annotators.
    \item \textbf{Rejection Sampling Fine-tuning (RFT)}: In this approach, the SFT model undergoes further fine-tuning using its own generated outputs for SFT prompts, where only those outputs passing a correctness-based filter are retained for further training.
    \item \textbf{Direct Preference Optimization (DPO)}: DPO builds upon SFT by further optimizing the model using an enhanced set of generated responses, applying a pairwise DPO loss to improve performance.
    \item \textbf{Online Rejection Sampling Fine-tuning (Online RFT)}: Distinct from conventional RFT, Online RFT leverages the SFT model to initialize the policy and iteratively fine-tunes it on outputs produced by the policy in real-time, adapting continuously with new samples.
    \item \textbf{PPO/GRPO}: These methods begin from the SFT policy, then apply reinforcement using outputs sampled directly from the current policy model during training.
\end{itemize}

The constituent elements of these approaches are presented in Table \ref{tab:data_policy}. For an in-depth explanation of the derivation, consult Appendix \ref{app:analysis-rl}.

\paragraph{Observation about Data Source} We categorize data sources into online and offline sampling. Online sampling involves generating training data from the ongoing exploratory actions of the current policy during training, whereas offline sampling relies on collecting training data from the sampling output of the pre-trained SFT model. Both RFT and DPO utilize offline data, while Online RFT and GRPO depend on data gathered online.

Figure \ref{fig:rl-analysis} illustrates that Online RFT achieves notably better performance than RFT across two benchmark tasks. In particular, at the onset of training, both approaches exhibit similar results; however, as training progresses, Online RFT establishes a clear lead, thereby highlighting the benefits of online sampling. This trend is expected: in early iterations, the actor and the SFT model are highly aligned, resulting in sampled data with minimal variation. As training advances, actor-generated samples begin to diverge more significantly, and leveraging online, real-time data becomes increasingly beneficial.

\paragraph{Observation about Gradient Coefficient} Reward signals are transformed into gradient coefficients that facilitate updates to model parameters. Our experimental setup includes two variants of reward functions: ‘Rule’ and ‘Model’. The ‘Rule’ approach evaluates response quality directly through answer correctness, while the ‘Model’ strategy utilizes a trained reward model—whose training is itself guided by rule-based assessments—to assign scores to each response. As indicated by Equations \ref{eq:GC-RFT} and \ref{eq:GC-GRPO}, GRPO introduces a distinctive mechanism by modulating the gradient coefficient according to the reward model’s output, thereby enabling varying strengths of reward or penalty contingent on response quality. By contrast, Online RFT does not distinguish among responses beyond correctness; incorrect answers are not penalized, and all correct responses are reinforced equally.

As illustrated in Figure \ref{fig:rl-analysis}, GRPO outperforms online RFT, underscoring the effectiveness of modifying both positive and negative gradient coefficients. Moreover, the results show that GRPO+PS achieves better outcomes than GRPO+OS, which demonstrates the advantage of incorporating step-specific, fine-grained gradient coefficients. We also investigate the effect of iterative RL by running two consecutive rounds in our experiments. According to Figure \ref{fig:iter-rl}, iterative RL yields notable performance gains, especially during the initial iteration.

\subsubsection{Why RL Works?} This work applies reinforcement learning to a subset of instruction tuning data, leading to substantial improvements over the baseline instruction-tuned model. To better understand the mechanism behind RL's success, we assess Pass@K and Maj@K accuracies for both the Instruct model and the RL-augmented model across two different benchmarks. Results depicted in Figure \ref{fig:maj-pass} reveal that RL primarily boosts the Maj@K metric, whereas Pass@K remains unchanged. These outcomes suggest that reinforcement learning predominantly strengthens the robustness of the model's output distribution—in effect, \textbf{the gains arise from an increased likelihood that the correct answer is ranked among the TopK outputs, rather than from advances in the model's underlying problem-solving ability.} Analogous observations have been made in prior work; for example, \citep{wang2023making} identified a \textbf{misalignment problem} in SFT models for reasoning tasks, and demonstrated that preference alignment strategies \citep{yuan2023rrhf,song2023preference,wang2023making} can meaningfully enhance SFT models’ reasoning performance.

\subsubsection{How to Achieve More Effective RL?}
\label{sec:effec_rl}
Our study confirms that reinforcement learning is highly effective for mathematical reasoning tasks. Furthermore, we introduce a unified framework to conceptualize a range of prominent training strategies. Under this framework, all approaches can be classified as either direct or approximate forms of RL. As expressed in Equation \ref{eq:objective}, three fundamental elements are present: Data Source, Algorithm, and Reward Function. We outline promising avenues for future research concerning each of these components.

\paragraph{Data Source} The data source provides the foundation for every training approach. Within the RL setting, this typically refers to unlabeled questions paired with outputs generated by the policy model. In this investigation, we utilize only questions from the instruction tuning stage and employ a simple nucleus sampling procedure for output generation. We hypothesize that this limitation contributes to the RL pipeline's effect being restricted to Maj@K improvement. Future work will consider expanding the RL framework to encompass out-of-distribution prompts and integrate \textbf{advanced sampling (decoding) strategies} such as tree-search-based methods \citep{yao2023tree}. Furthermore, the incorporation of \textbf{efficient inference techniques} \citep{xia-etal-2023-speculative,leviathan2023fast,kwon2023efficient,xia2024unlocking}, which are crucial for optimizing policy model exploration, represents another important research direction.

\paragraph{Algorithms} Algorithms translate the incoming data and reward feedback into gradient coefficients, which then drive updates to the model parameters. According to Equation \ref{eq:objective}, most current approaches inherently \textbf{TRUST} the signals from the reward function when adjusting the conditional probabilities assigned to specific tokens. Nevertheless, there is no guarantee that these reward signals maintain accuracy at all times, particularly in the context of highly intricate tasks. For instance, even datasets such as PRM800K \citep{lightman2023let}, which benefit from meticulous annotation by expert annotators, are reported to have around 20\% erroneous labels\footnote{\url{https://github.com/openai/prm800k/issues/12\#issuecomment-1728491852}}. Thus, we focus on investigating reinforcement learning algorithms designed for resilience against noisy or unreliable reward information. We posit that \textbf{WEAK-TO-STRONG} alignment strategies \citep{burns2023weak} have the potential to transform fundamental algorithmic paradigms.

\paragraph{Reward Function} The reward function provides the primary feedback signal during training. In reinforcement learning contexts, this typically takes the form of a neural reward model. We identify three crucial research directions for advancing reward models: 1) \textbf{Improving the generalization capabilities of reward models.} A robust reward model needs to generalize well to out-of-distribution inputs and novel decoding outputs; failure to do so may lead to reinforcement learning procedures that reinforce existing LLM distributions rather than genuinely enhancing core model abilities; 2) \textbf{Capturing and representing reward model uncertainty.} Accounting for uncertainty could serve as a valuable interface between weak reward models and weak-to-strong learning approaches; 3) \textbf{Developing efficient methodologies for high-fidelity process reward models} that deliver granular feedback during the reasoning phases \citep{lightman2023let,wang2023math}.

\section{Conclusion, Limitation, and Future Work} In this work, we introduce DeepSeekMath, which sets a new standard among open-source models on the competition-level MATH benchmark and achieves performance nearing that of proprietary systems. DeepSeekMath~was built by initializing from DeepSeek-Coder-v1.5 7B and conducting extended continual training over 500B tokens, incorporating a substantial subset of 120B math-specific tokens drawn from Common Crawl. Through thorough ablation studies, we demonstrate that mathematical data extracted from web sources holds considerable promise in quality, whereas data obtained from arXiv proves less beneficial than anticipated. Furthermore, we propose Group Relative Policy Optimization (GRPO), an adaptation of Proximal Policy Optimization (PPO), which enhances mathematical reasoning efficiency while reducing memory usage. Experimental evidence confirms the efficacy of GRPO, even when DeepSeekMath-Instruct 7B already achieves high benchmark scores. Additionally, we outline a unified framework for interpreting various methods, and discuss several promising avenues for future reinforcement learning improvements.

Despite DeepSeekMath's~strong results on benchmarks related to quantitative reasoning, the model's proficiency in geometry and theorem-proving is surpassed by leading closed models. For example, preliminary evaluations reveal that the model struggles with tasks involving triangles and ellipses, suggesting possible biases in the selection of training and fine-tuning data. Moreover, due to constraints in model size, DeepSeekMath~exhibits inferior few-shot learning abilities compared to GPT-4: while GPT-4's performance increases with few-shot prompts, DeepSeekMath~performs similarly in both zero-shot and few-shot scenarios. Moving forward, we intend to further refine our data selection methodologies to curate a superior pre-training corpus, and will pursue the research directions discussed in Section \ref{sec:effec_rl} to advance reinforcement learning techniques for LLMs.

\end{document}